% =====================================================================
%  MODELLO: DOCUMENTO  (Norme di Progetto, Piano di Progetto, Piano di
%  Qualifica, Analisi dei Requisiti, Glossario, Specifica Tecnica,
%  Manuale Utente, documenti della candidatura...)
%
%  1. Create una cartella per il documento e copiateci questo file:
%       <Milestone>/documenti/interni/<nome>/<nome>.tex   (uso interno)
%       <Milestone>/documenti/esterni/<nome>/<nome>.tex   (uso esterno)
%     es.  RTB/documenti/interni/norme-di-progetto/norme-di-progetto.tex
%     (Milestone = Candidatura, RTB oppure PB). Dopo il merge la CI
%     mette accanto il PDF con lo stesso nome; sul sito il nome ha in
%     fondo la versione. \uso qui sotto deve dire come la cartella.
%  2. Cambiate i campi qui sotto (dopo il simbolo % c'è un commento).
%  3. Scrivete il contenuto dopo \begin{document}.
%  Copertina, registro, indice, elenco di figure e tabelle (se ce ne
%  sono), intestazioni e numeri di pagina si fanno da soli.
%  Guida completa: GUIDA.md (ricetta 4).
% =====================================================================
\documentclass[documento]{gpt}

\titolo{Nome del documento}        % es. Norme di Progetto
\uso{esterno}                      % "interno" (es. Norme di Progetto)
                                   % oppure "esterno" (tutti gli altri)
\responsabile{spagnoli}            % chi approva il documento: il Responsabile

% Facoltativi (togliete il % davanti se servono):
% \capitolato{Cn -- Nome del capitolato (Azienda)}   % solo dopo l'assegnazione
% \destinatari{Prof.~Tullio Vardanega\newline Prof.~Riccardo Cardin\newline Azienda}

% ---------------------------------------------------------------------
%  REGISTRO DELLE MODIFICHE: qui non si scrive nulla.
%  Le righe le aggiunge il bot dopo il merge della pull request, nel file
%  nome-del-file.registro.tex accanto a questo, in base all'etichetta della PR:
%    contenuto     prima stesura e parti nuove   (0.1.0, 0.2.0, ...)
%    correzione    correzioni piccole             (0.2.1, ...)
%    approvazione  versione approvata             (1.0.0, ...)
%  Autore = chi ha fatto la PR, verificatore = chi l'ha approvata.
%  Approvazione (PR separata, etichetta messa dal Responsabile): nella riga
%  solo il Responsabile come autore, nessun verificatore.
%  stesura-e-approvazione (la mette il Responsabile): le due righe insieme,
%  contenuto (0.1.0) e approvazione (1.0.0), da una sola PR.
%  Regole: GUIDA.md, sezione 5.
% ---------------------------------------------------------------------

\begin{document}

\section{Introduzione}

\subsection{Scopo del documento}
A cosa serve questo documento e a chi si rivolge.

\subsection{Glossario}
I termini contrassegnati dalla lettera G a pedice, come \glo{repository},
sono spiegati nel Glossario.

\section{Primo argomento}
Testo.

% Documento lungo? Mettete ogni capitolo in un file a parte nella
% sottocartella capitoli/ e richiamatelo così (senza ".tex"):
% \input{capitoli/secondo-argomento}

\end{document}
